%%%%%%%%%%%%%%%%%%%%%%%%%
\spellentry{Baleful Polymorph}
%%%%%%%%%%%%%%%%%%%%%%%%%

\tagline{Transmutation}
\begin{description*}
\item[Level:] Drd 5, Sor/Wiz 5
\item[Casting Time:] 1 standard action
\item[Range:] Close (25 ft. + 5 ft./2 levels)
\item[Duration:] Permanent
\item[Saving Throw:] Fortitude negates, Will partial; see text
\item[Spell Resistance:] Yes
\item[Components:] V, S
\end{description*}

As \linkspell{Polymorph}, except that you change the subject into a Small or smaller
animal of no more than 1 HD. If the new form would prove fatal to the creature
the subject gets a +4 bonus on the save.
If the spell succeeds, the subject must also make a Will save. If this second save
fails, the creature loses its extraordinary, supernatural, and spell-like abilities,
loses its ability to cast spells (if it had the ability), and gains the alignment,
special abilities, and Intelligence, Wisdom, and Charisma scores of its new form
in place of its own. It still retains its class and level (or HD), as well as all
benefits deriving therefrom (such as base attack bonus, base save bonuses, and
hit points). It retains any class features (other than spellcasting) that aren't
extraordinary, supernatural, or spell-like abilities.
Incorporeal or gaseous creatures are immune to being polymorphed, and
a creature with the shapechanger subtype can revert to its natural form as a standard
action.